\documentclass[11pt, a4paper]{moderncv}

\moderncvstyle{classic}
\moderncvcolor{black}

\usepackage[scale=0.85, top=2cm, bottom=2cm]{geometry}

\usepackage[english]{babel}

\setlength{\hintscolumnwidth}{3cm}  

% ---------------------------------------------------------------
% Osobni informace
% ---------------------------------------------------------------
\name{Kamil}{Belán}
\title{Curriculum vitae}
%\born{22.4.2002}
%\address{Družstevní 20}{Třebíč}{Česká republika}
\phone[mobile]{+420~722~475~973}
\email{kamil@kamilbelan.net}
% \homepage{yourwebsite.com}
\social[github]{kamilbelan}
% \social[linkedin]{yourlinkedin}
% \photo[64pt][0pt]{photo}   % optional

% ---------------------------------------------------------------
\begin{document}

\makecvtitle

% ---------------------------------------------------------------
\section{Education}
\cventry{2026--now}{doctoral}{Faculty of Mathematics and Physics, Charles University (CUNI FMP)}{Prague}{Computational Mathematics}{ongoing}
\cventry{2024--2026}{master}{CUNI FMP}{Prague}{Mathematical and Computational Modelling in Physics}{}
\cventry{2021--2024}{bachelor}{CUNI FMP}{Prague}{Physics}{}
\cventry{2017--2021}{grammar school}{}{Třebíč}{}{}

% ---------------------------------------------------------------
\section{Research interests}
Numerical methods for partial differential equations of fluid dynamics; past work includes

\cvlistitem{stratified compressible atmospheric flow}
\cvlistitem{Lagrangian (particle) methods, in particular Smoothed Particle Hydrodynamics (SPH)}
\cvlistitem{well-balanced numerical schemes}

\smallskip

In the ongoing doctoral study the main research interests are higher order methods, in particular

\cvlistitem{discontinuous Galerkin method (DGM)}
\cvlistitem{$hp$-adaptive discontinuous Galerkin method}

% ---------------------------------------------------------------
\section{Theses}
\cventry{2026}{Numerical methods for the simulation of air flow past a mountain}{master's thesis}{supervisor doc. RNDr. Michal Pavelka, Ph.D.}{}{Objective of the theses is to further develop and verify the SPH method for compressible atmospheric flow around orography. The work relies on a reformulation of the problem in natural variables, robust terrain-following orography description and improved boundary conditions treatment at the inflow and outflow. Hydrostatic (in)stability is addressed by the development of novel well-balanced schemes and adaptive techniques. }
\cvitem{Codebase}{(both theses) \url{https://github.com/kamilbelan/sph-mountain-waves}}
\cventry{2024}{Simulating flows past a mountain range using smoothed particle hydrodynamics}{bachelor's thesis}{supervisor doc. RNDr. Michal Pavelka, Ph.D.}{}{The main objective of the thesis is to use the Lagrangian numerical method SPH to simulate compressible atmospheric flow over mountain orography, a problem typically solved by mesh methods (FVM, FD). It was shown by numerical evidence that the simulation of hydrostatic balance of a stratified atmosphere by the standard SPH method is so unstable, that SPH in its standard formulation is completely unable to capture phenomena like internal gravity waves.}

% ---------------------------------------------------------------
%\section{Projekty}

% ---------------------------------------------------------------
%\section{Výuka}

% ---------------------------------------------------------------
\section{Technical skills}
\cvitem{Programming languages}{Julia, Python, Mathematica, MATLAB}
\cvitem{Tools}{LaTeX, Git, SSH, bash, SLURM, Docker}

\section{Languages}
\cvitem{Czech}{native}
\cvitem{English}{\textbf{C1} -- Cambridge certificate}
\cvitem{German}{\textbf{\textasciitilde B1} -- courses at Charles University}

%\section{Pedagogické zkušenosti a certifikáty}


%\section{Další zkušenosti}


% ---------------------------------------------------------------
%\section{Ocenění a granty}

% ---------------------------------------------------------------

\end{document}
